%% ============================================================================
%% template.tex  --  Literature Review: the single master template
%%
%% Copy this file to start a review, keep lrreview.cls and fonts/ beside it.
%% Build with XeLaTeX, twice:   xelatex template.tex   (or: latexmk -xelatex)
%% Style rules live in lrreview.cls only. Never add margins, spacing or font
%% sizes in a report. If a rule is missing, add it to lrreview.cls.
%% No manual page breaks. No em dashes, hash signs or asterisks in the text.
%% Citations in the text are author-year, with no hyperlinks and no numbers.
%% Replace every bracketed placeholder below.
%% ============================================================================
\documentclass{lrreview}

%% ---- Metadata (front page and running headers) -----------------------------
\lrtitle{[Exact paper title as published]}
\lrrunningtitle{[Short running title]}
\lrauthors{[Author One, Author Two and Author Three]}
\lryear{[Year]}
\lrjournal{[Journal, volume(issue), pages, or working paper series]}
\lrdoi{[DOI]}
\lrarea{[Research area and subfield]}
\lrseries{[Series name exactly as in CLAUDE.md]}
\lrpapertype{[Theoretical, Empirical, Econometric, Computational, Experimental, Review, Methodological or Mixed]}
\lrreviewer{Arthur Faugeron}
\lrdatereviewed{[D Month YYYY]}

\begin{document}

%% ---- Front page: metadata table and one analytical abstract (150 to 220 words)
\lrfrontpage{%
[One tight paragraph of about 150 to 220 words in the reviewer's own words: the
question, the data or model, the method, the main finding and the main limitation.
It must end as a complete sentence so that it fits on the front page.]}

%% ============================================================================
\section{Introduction}

\subsection{Paper identification and classification}
[Paper type, field and subfield, version of the paper read in full.]

\subsection{Executive understanding}
\subsubsection{What is the paper trying to discover?}
[Text.]
\subsubsection{Why does the question matter?}
[Text.]
\subsubsection{What was known before, and what is missing?}
[Text.]
\subsubsection{What does the paper propose, how does it investigate, and what does it find?}
[Text.]

%% ============================================================================
\section{Research Problem and Hypothesis}

\subsection{Research question}
\lrtag{Paper} [The precise question, not a topic description.]

\subsection{Hypotheses}
\begin{itemize}
  \item\relax [First hypothesis.]
  \item\relax [Second hypothesis.]
\end{itemize}

\subsection{Existing literature and research gap}
[Text. Use the labels Paper, Evidence, Critique, Inference and Extension at the
start of a paragraph so that the source of each statement is never ambiguous.]

\subsection{Claimed contribution}
\lrtag{Critique} [Whether the claimed contribution is actually demonstrated.]

%% ============================================================================
\section{Theoretical Framework and Mechanism}

\subsection{[Model or mechanism]}
\lrtag{Paper} [Every important equation must be shown, not only described.]
\begin{equation}
  y_t = \alpha + \beta\, x_t + \varepsilon_t .
  \label{eq:main}
\end{equation}
\begin{lrdefs}
  \lrdef{$y_t$}{[Meaning of each symbol.]}
  \lrdef{$x_t$}{[Meaning.]}
  \lrdef{$\beta$}{[Meaning, sign and units.]}
\end{lrdefs}
[What the equation says, which assumptions it needs, what it predicts, how the
prediction is tested, and the derivation steps. Refer to it as Equation
(\ref{eq:main}).]

%% ============================================================================
\section{Data and Empirical Strategy}

\subsection{Data}
[Source, sample period, frequency, geography, observations, filters, transformations.]

\begin{lrtable}{[Table caption]}\label{tab:data}
\begin{tabularx}{\linewidth}{@{}L L L r@{}}
\lrtoprule
\lrth{Series} & \lrth{Construction} & \lrth{Sample} & \lrth{Observations}\\
\lrheadrule
\relax [Series one] & [Construction] & [Period] & [N]\\
\lrrowrule
\relax [Series two] & [Construction] & [Period] & [N]\\
\lrbottomrule
\end{tabularx}
\lrsource{Source: [Origin of the numbers; say whether copied from the paper (official) or computed by the reviewer (generated).]}
\end{lrtable}

\subsection{Empirical specification and identification}
[Where does the identifying variation come from?]

\subsection{Data limitations and biases}
\begin{itemize}
  \item\relax [Survivorship, look-ahead, selection, measurement error, data mining.]
\end{itemize}

%% ============================================================================
\section{Results}

\subsection{[Main result]}
\lrtag{Evidence} [What was tested, found, its magnitude, statistical and economic
significance, and what it means in intuitive units.]

%% Figure example (needs an image file; keep it commented until one exists):
%% \lrfigure{fig_example.png}
%%   {[Figure title with units and period]}
%%   {Source: [official source or generated by the reviewer]; [sample, units, notes].}

%% ============================================================================
\section{Robustness and Additional Tests}

\begin{lrtable}{Robustness tests: concern, test, result and assessment}\label{tab:rob}
\lrsans\fontsize{8.4}{10.5}\selectfont
\begin{xltabular}{\linewidth}{@{}L L L L@{}}
\lrtoprule
\lrth{Concern} & \lrth{Test reported} & \lrth{Result} & \lrth{Does it resolve the concern?}\\
\lrheadrule
\endhead
\lrbottomrule
\endfoot
\relax [Concern] & [Test] & [Result] & [Yes, Partly or No, and why]\\
\lrrowrule
\relax [Concern] & [Test] & [Result] & [Assessment]\\
\end{xltabular}
\lrsource{Source: [Origin].}
\end{lrtable}

%% ============================================================================
\section{Critical Evaluation}

\subsection{What the paper establishes}
\begin{lrfinding}{Established}
[Claims strongly supported by the evidence.]
\end{lrfinding}
\begin{lrfinding}{Supported but conditional}
[Claims that depend on assumptions or particular specifications.]
\end{lrfinding}
\begin{lrfinding}{Suggestive}
[Evidence consistent with a hypothesis but insufficient to establish it.]
\end{lrfinding}
\begin{lrfinding}{Not established}
[Claims the paper cannot demonstrate.]
\end{lrfinding}

\subsection{Strengths}
\subsubsection{Strengths recognised by the authors}
\begin{itemize}
  \item\relax [Strength and why it matters.]
\end{itemize}
\subsubsection{Strengths identified independently}
\begin{itemize}
  \item\relax [Strength and why it matters.]
\end{itemize}

\subsection{Weaknesses and critique}
\subsubsection{Single most important weakness}
\textbf{Problem.} [Text.] \textbf{Why it matters.} [Text.] \textbf{Potential bias.} [Text.]
\textbf{How it could be tested.} [Text.]
\subsubsection{Other weaknesses}
\begin{itemize}
  \item\relax [Problem, why it matters, potential bias, test.]
\end{itemize}

%% ============================================================================
\section{Contradictions and Edge Cases}
\lrtag{Critique} [Conditions under which the result should disappear, extreme
values, parameter dependence, alternative mechanisms, with direct reasoning or
illustrative calculation.]

%% ============================================================================
\section{Further Literature}

\begin{lrtable}{Literature connections}\label{tab:lit}
\lrsans\fontsize{8.4}{10.5}\selectfont
\begin{xltabular}{\linewidth}{@{}L L L L@{}}
\lrtoprule
\lrth{Paper} & \lrth{Shared idea} & \lrth{Relationship} & \lrth{Why it matters}\\
\lrheadrule
\endhead
\lrbottomrule
\endfoot
\relax [Author (Year)] & [Idea] & [Preceded, contradicted, extended, applied] & [Reason]\\
\lrrowrule
\relax [Author (Year)] & [Idea] & [Relationship] & [Reason]\\
\end{xltabular}
\lrsource{Source: [Origin].}
\end{lrtable}

%% ============================================================================
\section{Research Gaps}

\subsection{Paper-level gaps}
\begin{itemize}
  \item \textbf{Gap:} [Text.] \emph{Why it matters:} [Text.] \emph{Existing limitation:} [Text.] \emph{Way forward:} [Text.]
\end{itemize}
\subsection{Literature-level gaps}
\begin{itemize}
  \item\relax [Same four elements.]
\end{itemize}
\subsection{Methodological gaps}
\begin{itemize}
  \item\relax [Same four elements.]
\end{itemize}

%% ============================================================================
\section{Research Opportunities}

\begin{lridea}{Idea 1. [Title]}
\lrfield{Research question}{[Text.]}
\lrfield{Motivation}{[Text.]}
\lrfield{Hypothesis}{[Text.]}
\lrfield{Economic mechanism}{[Text.]}
\lrfield{Literature connection}{[Text.]}
\lrfield{Data}{[Text.]}
\lrfield{Variables}{[Text.]}
\lrfield{Method}{[Text.]}
\lrfield{Identification strategy}{[Text.]}
\lrfield{Expected contribution}{[Text.]}
\lrfield{Main risk}{[Text.]}
\end{lridea}

%% Repeat the lridea block for about five ideas.

\subsection{Ranking}
\Needspace*{13\baselineskip}
\begin{lrtable}{Ranking of research ideas (scores from 1 to 5, assigned by the reviewer)}
\begin{tabularx}{\linewidth}{@{}L r r r r r r r@{}}
\lrtoprule
\lrth{Idea} & \lrth{Novelty} & \lrth{Feasibility} & \lrth{Data} & \lrth{Method} & \lrth{Importance} & \lrth{Contribution} & \lrth{Total}\\
\lrheadrule
\relax [1. Idea] & [n] & [n] & [n] & [n] & [n] & [n] & [n]\\
\lrbottomrule
\end{tabularx}
\lrsource{Source: reviewer's subjective scoring.}
\end{lrtable}
[The top three ideas and why.]

%% ============================================================================
\section{Learning Framework}

\subsection{Prerequisite concepts}
\begin{itemize}
  \item\relax [Concept.]
\end{itemize}
\subsection{Core concepts learned}
\begin{itemize}
  \item\relax [Concept.]
\end{itemize}
\subsection{Important equations}
[Equations worth remembering, referred to by number.]
\subsection{Methods learned}
[Text.]
\subsection{Questions to test understanding}
\subsubsection{Comprehension}
\begin{enumerate}
  \item\relax [Question.]
\end{enumerate}
\subsubsection{Analysis}
\begin{enumerate}
  \item\relax [Question.]
\end{enumerate}
\subsubsection{Research}
\begin{enumerate}
  \item\relax [Question.]
\end{enumerate}
\subsection{Papers to read next}
[Authors (Year); Authors (Year).]

%% ============================================================================
\section{Conclusion}

\subsection{One-sentence thesis}
[Text.]
\subsection{Three most important insights}
\begin{itemize}
  \item\relax [Insight.]
\end{itemize}
\subsection{Strongest evidence}
[Text.]
\subsection{Biggest weakness}
[Text.]
\subsection{Most important research gap}
[Text.]
\subsection{Best research opportunity}
[Text.]
\subsection{Three papers to read next}
[Authors (Year); Authors (Year); Authors (Year).]

%% ============================================================================
\section{References}
\begin{lrbibliography}
\lrref{[Author, A., and Author, B. (Year). Title. \emph{Journal}, volume(issue), pages. Add the version read, for example arXiv identifier and version.]}
\lrref{[Next reference, alphabetical by first author.]}
\end{lrbibliography}

\end{document}
